\documentclass[12pt]{article}
\usepackage{geometry}
\usepackage{amsfonts}
\usepackage{float}
\usepackage{amsmath}

\geometry{
legalpaper, total={177.8mm, 290mm},left=5mm, right=15mm,
top=7mm, bottom=12mm,
}

\usepackage{letltxmacro}

\LetLtxMacro\Oldfootnote\footnote

\newcommand{\EnableFootNotes}{%
  \LetLtxMacro\footnote\Oldfootnote%
}

\newcommand{\DisableFootNotes}{%
  \renewcommand{\footnote}[2][]{\relax}
}

\DisableFootNotes

\begin{document}

\iffalse
\begin{table}[]
\begin{tabular}{llcllcll}
 &  & \textbf{MYMENSINGH GIRLS’  CADET COLLEGE} &  &                                            & \multicolumn{1}{l}{}            &                        &                        \\
 &  & WHAT EXAMINATION - 2025                   &  &                                            &                                 &                        &                        \\
 &  & CLASS: WHAT                               &  &                                            &                                 &                        &                        \\
 &  & STATISTICS (CREATIVE)                     &  &                                            &                                 &                        &                        \\
 &  & WHAT PAPER                                &  &                                            & \multicolumn{1}{r}{}            &                        &                        \\ \cline{7-8} 
 &  & [According to the Syllabus of 2025]       &  & \multicolumn{1}{r}{\textbf{Subject Code:}} & \multicolumn{1}{l|}{\textbf{1}} & \multicolumn{1}{l|}{3} & \multicolumn{1}{l|}{0} \\ \cline{7-8} 
 &  & TIME – 2 hours \& 35 minutes              &  &                                            & \multicolumn{1}{r}{}            &                        &                        \\
 &  & FULL MARKS – 50                           &  &                                            & \multicolumn{1}{r}{\textbf{}}   &                        &                       
\end{tabular}
\end{table}
\fi

\begin{table}[]
\begin{tabular}{llllllllcllcl}
\textit{Proxima} &  &  &  &  &  &  &  & \textbf{MYMENSINGH GIRLS’ CADET COLLEGE} &                                             &                                 & \multicolumn{1}{l}{}            &                        \\
                 &  &  &  &  &  &  &  & YEAR FINAL EXAMINATION - 2026          &                                             &                                 &                                 &                        \\
                 &  &  &  &  &  &  &  & CLASS: XI                               &                                             &                                 &                                 &                        \\
                 &  &  &  &  &  &  &  & STATISTICS (CREATIVE)                    &                                             &                                 &                                 &                        \\
                 &  &  &  &  &  &  &  & FIRST PAPER                             &                                             &                                 & \multicolumn{1}{r}{}            &                        \\ \cline{11-13} 
                 &  &  &  &  &  &  &  & [According to the Syllabus of 2027]      & \multicolumn{1}{r|}{\textbf{Subject Code:}} & \multicolumn{1}{l|}{\textbf{1}} & \multicolumn{1}{l|}{\textbf{2}} & \multicolumn{1}{l|}{9} \\ \cline{11-13} 
                 &  &  &  &  &  &  &  & TIME – 2 hours \& 35 minutes             &                                             &                                 & \multicolumn{1}{r}{}            &                        \\
                 &  &  &  &  &  &  &  & FULL MARKS – 50                          &                       
\end{tabular}
\end{table}

\vspace{-1cm}
\hrule

\begin{center}
[\textbf{N.B.} – The figures of the right margin indicate full marks. Read the stems carefully and answer the associated questions. Answer any \textbf{FIVE} questions taking at least two questions from each group]\\
\end{center}

\begin{center}
\textbf{Group  - A}
\end{center}
\begin{enumerate}

\item

  \textbf{Number of accidents occurred in two different cities A and B are denoted by two variables x and y, respectively.} 
  
  \begin{table}[h]
  \centering
\begin{tabular}{c|c|c|c}
City & Day 1 & Day 2 & Day 3 \\ \hline
A    & 3     & 4     & 2     \\
B    & 5     & 3     & 2    
\end{tabular}
\end{table}
  
  An analyst observed that $\displaystyle \Bigl(\sum_{i=1}^{3} x_i\Bigr)^2 \;\neq\; \sum_{i=1}^{3} x_i^2$ and  $\displaystyle \Bigl(\sum_{i=1}^{3} y_i\Bigr)^2 \;\neq\; \sum_{i=1}^{3} y_i^2$
  
  \begin{enumerate}
   \item What is a quantitative variable? \hfill 1
 \item Distinguish qualitative and quantitative variables. \hfill 2
    \item  Show that 
	\[
\sum_{i=1}^{3} (x_i + y_i)^2 = \sum_{i=1}^{3} x_i^2 + 2\sum_{i=1}^{3} x_i y_i + \sum_{i=1}^{3} y_i^2
\] \hfill 3
    \item
	Evaluate the analyst's comment. \hfill 4
  \end{enumerate}
  
    \item
\textbf{The test scores (out of 100) of 20 students in a mathematics exam were recorded as follows:}
\begin{center}
65, 72, 78, 81, 55, 88, 92, 47, 69, 73 \\
84, 91, 77, 63, 58, 71, 94, 86, 68, 79 
\end{center}

\begin{enumerate}
    \item
    What are discrete class intervals? \hfill 1
    \item
    What are the limitations of using discrete class intervals with continuous data? \hfill 2
    \item
    Make a frequency distribution and explain. \hfill 3
    \item  
    Construct a histogram for the given data and interpret its shape. \hfill 4
\end{enumerate}

   \item 
  \textbf{Two schools, X and Y, conducted the same mathematics examination for their students. The summary of the test scores (out of 100) is shown below:}

\begin{table}[h]
\centering
\begin{tabular}{c|ccc}
School & Mean Score & Standard Deviation & Number of Students \\ \hline
X      & 72.5       & 8.2                & 150                \\ 
Y      & 74.0       & 6.5                & 120               
\end{tabular}
\end{table}

  \begin{enumerate}
      \item When is variance less than standard deviation? \hfill 1
      \item Is coefficient of variation a pure number? Explain \hfill 2
    \item  
	Determine the combined standard deviation of the scores of the schools. \hfill 3
    \item
	Which school has more consistent scores? Analyze. \hfill 4
  \end{enumerate}

  \begin{center}
\textbf{Group  - B}
\end{center}


  
   \item
	  \textbf{The first four moments around 4 of productions of a company over four years are -1.5, 17, -30, and 108.} 
  
  \begin{enumerate}
  \item  What is raw moment? \hfill 1
    \item
	Can the second central moment be negative? \hfill 2
    \item  
	Determine the second and third central moments. \hfill 3
    \item
	What kind of kurtosis do the data have? \hfill 4
  \end{enumerate}

   \item
	  \textbf{The following table shows the exam scores of 10 students and the number of hours they studied for the exam. It is hypothesized that the exam score depends on the number of hours studied.} 
	  
\begin{table}[h!]
\centering
\begin{tabular}{c|c|c|c|c|c|c|c|c|c|c}
Hours Studied (X) & 2 & 3 & 4 & 5 & 6 & 7 & 8 & 9 & 10 & 11 \\ \hline
Exam Score (Y) & 65 & 70 & 72 & 78 & 80 & 85 & 88 & 92 & 95 & 98 \\
\end{tabular}
\end{table}

  \begin{enumerate}
    \item What is a scatter plot? \hfill 1
  \item Differentiate between the values of correlation $r_1 = +1$ and $r_2=-1$ \hfill 2
    \item  
	Make a scatter plot using the data and state the potential direction and strength of the relationship between the variables. \hfill 3
    \item
	Find the correlation and interpret. \hfill 4
  \end{enumerate}

  \item
\textbf{The monthly sales (in units) of laptops at a major electronics store are recorded below:} 

\begin{table}[h]
\centering
\begin{tabular}{c|cccccccc}
\textbf{Month} & Jan & Feb & Mar & Apr & May & Jun & Jul & Aug \\ \hline
\textbf{Sales} & 120 & 150 & 130 & 170 & 160 & 180 & 200 & 210
\end{tabular}
\end{table}

\begin{enumerate}
  \item What is an irregular variation? \hfill 1
   \item What are the components of Time Series?   \hfill 2
  \item
  Illustrate the trend graphically and estimate the expected laptop sales for September. \hfill 3
  \item  
  Compute the trend using a three-monthly moving average method and exraplotae to September, and compare with the graphical method. \hfill 4

\end{enumerate}

\end{enumerate}
 \vspace{2.5cm}

\begin{center}

  \vfill
  --AAM--
\end{center}



\end{document}